% Generated draft: requires mathematical and source review.
\documentclass[11pt]{article}
\usepackage[margin=25mm]{geometry}
\usepackage{fontspec}
\setmainfont{Latin Modern Roman}
\usepackage{amsmath,amssymb,mathtools,unicode-math}
\setmathfont{Latin Modern Math}
\usepackage{enumitem,needspace,xurl,hyperref}
\setlength{\parindent}{0pt}
\setlength{\parskip}{5pt}
\setlength{\emergencystretch}{3em}
\newcommand{\N}{\mathbb{N}}
\newcommand{\Z}{\mathbb{Z}}
\newcommand{\Q}{\mathbb{Q}}
\newcommand{\R}{\mathbb{R}}
\newcommand{\C}{\mathbb{C}}
\newcommand{\F}{\mathbb{F}}
\newcommand{\A}{\mathbb{A}}
\newcommand{\PP}{\mathbb{P}}
\newcommand{\E}{\mathbb{E}}
\newcommand{\eps}{\varepsilon}
\newcommand{\dd}{\,\mathrm{d}}
\begin{document}
\section*{UnsolvedMath ID: 30004586}
\textbf{Source problem code:} catalogue entry
\subsection{Definitions and mathematical statement}
\textbf{Scope and unresolved conventions.} The supplied record asks about the possible spectra of weighted Laplace or Markov operators on graphs arising from groups, with two particular existence questions. It does not specify the admissible graphs, weights, operator domains, Hilbert spaces, or meaning of a spectral gap. Consequently, the following standard specialization makes the questions explicit but is not asserted to be the uniquely intended formulation.

Let \(G\) be a finitely generated group, let \(S\subseteq G\setminus\{e\}\) be a finite symmetric generating set, and let \(p:S\to(0,1]\) satisfy
\[
 p(s)=p(s^{-1}),\qquad \sum_{s\in S}p(s)=1.
\]
The undirected Cayley graph \(\operatorname{Cay}(G,S)\) has vertex set \(G\), with an edge between \(g\) and \(gs\) for each \(g\in G\) and \(s\in S\). On \(\ell^2(G)\) with counting measure, define the symmetric Markov operator and normalized Laplace operator by
\[
 (P_pf)(g)=\sum_{s\in S}p(s)f(gs),\qquad L_p=I-P_p.
\]
These are bounded self-adjoint operators. For a bounded operator \(T\), write
\[
 \sigma(T)=\{\lambda\in\C:T-\lambda I\text{ has no bounded inverse}\}.
\]
Here \(\sigma(P_p)\subseteq[-1,1]\), \(\sigma(L_p)\subseteq[0,2]\), and \(\sigma(L_p)=\{1-\lambda:\lambda\in\sigma(P_p)\}\). A Cantor set means a nonempty compact subset of \(\R\) with no isolated points and empty interior. A group \(G\) is torsion-free if \(g^n=e\), for \(g\in G\) and an integer \(n\geq1\), implies \(g=e\).

\textbf{Questions from the supplied record.} Characterize the subsets of the spectral plane that can occur as spectra of weighted Laplace or Markov operators on graphs arising from groups, once the admissible graph and operator class has been specified. In the symmetric Cayley specialization above, the Cantor-spectrum existence question is
\[
 \exists\,(G,S,p)\quad \sigma(P_p)\text{ is a Cantor set},
\]
 equivalently with \(L_p\) in place of \(P_p\). The source does not determine whether arbitrary symmetric weights are allowed, whether only the uniform weights \(p(s)=|S|^{-1}\) are intended, or whether this specialization captures its full target.

The second existence question asks whether a Cayley graph of a torsion-free group can have a spectral gap. This cannot be given a unique quantified formulation from the supplied record. For example, in the specialization above, possible meanings include a positive lower spectral bound
\[
 \exists\,\eps>0\quad \inf\sigma(L_p)\geq\eps,
\]
 a gap above a spectral point at zero
\[
 0\in\sigma(L_p),\qquad
 \exists\,\eps>0\quad \sigma(L_p)\cap(0,\eps)=\varnothing,
\]
 or an internal gap
\[
 \exists\,a,b\in\sigma(L_p)\quad
 a<b,\qquad (a,b)\cap\sigma(L_p)=\varnothing.
\]
For each interpretation, the intended existence quantifier would range over admissible triples \((G,S,p)\) with \(G\) torsion-free. These interpretations are different; the record does not select one. It also does not specify whether \(G\) must be infinite or whether the operators must be self-adjoint. These unresolved qualifications must be settled before a definitive mathematical formulation is accepted.
\subsection{Short English statement}
The problem asks which spectral shapes occur for weighted Laplace or Markov operators on graphs associated with groups, whether a Cayley graph can have Cantor spectrum, and whether one associated with a torsion-free group can have a spectral gap. A standard symmetric Cayley-graph model is stated explicitly, but the intended weights, operator setting, graph class, and meaning of spectral gap remain unspecified. The supplied open status is dated curation, not an independently verified conclusion.
\subsection{Sources}
\begin{itemize}
\item Supplied local UnsolvedMath registry record, permanent ID \(30004586\), title \emph{Spectral Shapes of Group-Based Graph Operators}, source problem code \texttt{catalogue entry}, in the \emph{Oberwolfach Reports Open Problems} collection. No individual proposer is supplied; the proposed year is recorded as 2021.
\item Supplied report citation: \emph{Geometry, Dynamics and Spectrum of Operators on Discrete Spaces} (2022), \emph{Oberwolfach Reports}, DOI 10.4171/owr/2021/2, pp.\ 33--85. \url{https://doi.org/10.4171/owr/2021/2}. This is the report attribution supplied by the registry; its bibliographic details and precise problem location have not been independently verified.
\item Supplied literature source, labeled ``spectral graph survey'': \url{https://access.archive-ouverte.unige.ch/access/metadata/2e554b7d-359c-4c36-92ff-fd452566e950/download}. The record supplies no author, title, or precise passage for this source.
\item Supplied literature assessment checked on \textbf{2026-08-22}: the record labels the problem open and reports that the cited survey identifies both Cantor spectrum for Cayley graphs and spectral gaps for torsion-free groups as major open problems, with later literature continuing to describe the Cayley-graph question as open. This is supplied dated curation, not an independently verified literature assessment. The background explicitly requests independent verification of time-sensitive status and cited claims.
\end{itemize}
\end{document}
